\section{Objetivos y preguntas}

\subsection{Preguntas de Investigación}
\begin{itemize}
    \item \textbf{Pregunta General:} ¿Cuáles son los mecanismos organizacionales, condiciones necesarias y sesgos de género que explican la adopción empresarial y la permanencia laboral de mujeres con antecedentes penales en la red de Juntos por la Reinserción (JxR)?
    \item \textbf{Preguntas Específicas:}
    \begin{enumerate}
        \item ¿Cómo se comportan las tasas de permanencia e interrupción laboral según género en los registros históricos y administrativos de JxR?
        \item ¿De qué manera condiciona el género de la candidatura la disposición de contratación bajo perfiles equivalentes (evaluado mediante viñetas pareadas)?
        \item ¿Qué configuraciones de condiciones organizacionales (QCA) resultan necesarias o suficientes para la adopción empresarial?
        \item ¿Qué procesos causales y puntos de quiebre explican la permanencia versus la interrupción laboral en las relaciones empresa--trabajadora?
        \item ¿Qué ajustes operacionales requieren los tres modelos de JxR (Bolsa Laboral, OTEC y CET) según el tipo de empresa?
    \end{enumerate}
\end{itemize}

\subsection{Objetivos del Estudio}
\begin{itemize}
    \item \textbf{Objetivo General:} Explicar mediante un estudio de casos múltiples los mecanismos de adopción organizacional y permanencia laboral de mujeres en proceso de reinserción, proponiendo una tipología B2B de empresas y el rediseño programático de la oferta de JxR.
    \item \textbf{Objetivos Específicos:}
    \begin{enumerate}
        \item Establecer la línea base de permanencia e interrupción laboral por género a partir del \textit{desk research} de registros operativos de JxR y datos sectoriales.
        \item Evaluar el efecto del género en la disposición B2B de contratación mediante viñetas pareadas intra-sujeto.
        \item Modelar las combinaciones de condiciones causales de adopción empresarial mediante Análisis Cualitativo Comparado (QCA).
        \item Identificar los mecanismos de permanencia y quiebre mediante comparación emparejada y rastreo de procesos en díadas empresa--trabajadora.
        \item Formular una tipología B2B de empresas y adaptar los modelos de colaboración (Bolsa Laboral, OTEC y CET) según los requerimientos detectados.
    \end{enumerate}
\end{itemize}